\section{Methodology}
\label{sec:method}

\subsection{Profile Coherence}
\label{sec:method:pck}

\subsubsection{Ordinal band scale}
Customers and assets share the same 4-point ordinal scale running from Conservative to Income to Balanced to Aggressive, encoded 0 to 3. The customer band $b_u$ is taken from FAR-Trans's \texttt{riskLevel} field. FAR-Trans provides both declared values and regression-imputed \texttt{Predicted\_*} entries, and both are mapped to the same ordinal scale. Customers whose entry is \texttt{Not\_Available} or missing are excluded from all coherence computations. The asset band $b_i$ follows the hierarchical rule in Section~\ref{sec:setup:bandassignment}.

\subsubsection{Pairwise discordance}
\begin{equation}
d(u, i) = |b_u - b_i| \in \{0, 1, 2, 3\}.
\end{equation}
A recommendation is \emph{profile-coherent} iff $d \leq 1$. The tolerance of $\pm 1$ band reflects MiFID II's adjacency allowance for borderline suitability. Section~\ref{sec:limitations} discusses this design choice further.

\subsubsection{Aggregation}
\begin{equation}
\text{PC@}k(u) = \frac{1}{k} \left|\{ i \in \text{top}_k(u) : d(u, i) \leq 1 \}\right|.
\end{equation}
Two implementation rules complete the definition. First, recommenders returning fewer than $k$ items use the actual number of items returned as the denominator. Second, customers with $b_u = \emptyset$ are dropped from the aggregate. Per-split PC@$k$ is the unweighted mean across eligible customers and the headline numerics average over all splits.

\subsubsection{Band-conditional random baseline}
The closed-form expectation of PC@$k$ under a uniformly random recommender for band $b$ is given by
\begin{equation}
\pi(b) = \frac{|\{ i \in A : |b - b_i| \leq 1 \}|}{|A|},
\end{equation}
where $A$ is the asset universe. Section~\ref{sec:setup:bandassignment} reports the FAR-Trans asset-band distribution and the resulting $\pi(b)$ values.

\subsubsection{PC-lift@$k$}
Per-user PC@$k$ is normalised against the random reference at the customer's own band:
\begin{equation}
\text{PC-lift@}k(u) = \frac{\text{PC@}k(u)}{\pi(b_u)}.
\end{equation}
Lift $> 1$ beats the band-conditional random rate, while lift $= 1$ matches it. Because 77.9\% of FAR-Trans assets sit in the Conservative, Income, and Balanced bands, a recommender that pushes Income assets at everyone would score well on raw PC@$k$ for any user whose tolerance set includes Income. Dividing by $\pi(b_u)$ strips the band-specific easiness of the asset universe and makes per-band values commensurate.

\subsection{Stratified, Profile-Coherent LightGCN}
\label{sec:method:pclgcn}

\subsubsection{Backbone}
We choose LightGCN as the backbone because its differentiable loss and per-customer embeddings allow us to attach a coherence margin directly to the BPR objective. Random Forest is not extended because it ranks assets by predicted return from technical indicators alone and carries no customer-level signal to condition on.

\subsubsection{Stratified architecture}
We train four band-stratified sub-models, one per declared MiFID risk band (Figure~\ref{fig:architecture}). Each sub-model's training graph contains only the interactions of customers assigned to that band. At inference, each customer is routed to the sub-model matching their declared band. Customers with no declared band default to the Balanced sub-model.

\subsubsection{Profile-coherent margin loss}
BPR trains the model to score an asset the customer actually bought (the positive item) higher than a randomly sampled unobserved asset (the negative item). Independently of this BPR sampling, we add a second margin term that samples one coherent asset $i_c$ (satisfying $d(u, i_c) \leq 1$) and one discordant asset $i_d$ (satisfying $d(u, i_d) > 1$) from the sub-model's asset universe, and pushes the coherent score above the discordant score:
\begin{equation}
\mathcal{L}_{\text{coh}}(u) = \text{softplus}\!\left(s(u, i_d) - s(u, i_c)\right),
\end{equation}
where $s(u, i)$ is the LightGCN inner-product score between customer $u$ and asset $i$.

\subsubsection{Total loss}
\begin{equation}
\mathcal{L} = \mathcal{L}_{\text{BPR}} + \omega \|\Theta\|_2^2 + \lambda \cdot \mathcal{L}_{\text{coh}},
\end{equation}
Here $\omega$ is the L2 weight decay on embedding parameters $\Theta$ and $\lambda \geq 0$ controls how strongly the coherence term influences training. BPR remains the primary ranking signal, and only the treatment trial activates the coherence term.

\subsubsection{Ablation Setup}
Setting $\lambda = 0$ trains stratified BPR without the coherence term, isolating the effect of the stratified architecture. Setting $\lambda = 1$ adds the coherence margin on top. Comparing the two trials lets us attribute gains separately to stratification and to the PC-loss.

\begin{figure*}[!t]
\centering
\resizebox{0.85\textwidth}{!}{%
\begin{tikzpicture}[
    >=Stealth,
    node distance=0.35cm and 0.5cm,
    box/.style={draw, rounded corners=2pt, minimum height=0.5cm,
                minimum width=1.2cm, align=center, font=\scriptsize},
    lossbox/.style={draw, rounded corners=2pt, minimum height=0.45cm,
                    minimum width=0.9cm, align=center, font=\scriptsize,
                    fill=gray!15},
    groupbox/.style={draw, dashed, rounded corners=4pt, inner sep=3pt},
    arr/.style={->, thick},
    line/.style={thick},
]

\node[box] (cust) {Customer $u$\\band $b_u$};

\node[box, right=0.5cm of cust] (router) {Band\\Router};

\coordinate (subanchor) at ($(router.east) + (1.3,0.9)$);
\node[box, anchor=west] at (subanchor) (sub0) {Cons.};
\node[box, below=0.1cm of sub0] (sub1) {Income};
\node[box, below=0.1cm of sub1] (sub2) {Balanced};
\node[box, below=0.1cm of sub2] (sub3) {Aggr.};

\node[box, right=2.0cm of $(sub1)!0.5!(sub2)$] (scores) {Scores\\$s(u, i)$};

\node[box, right=0.6cm of scores] (topk) {Top-$k$\\list};

\node[lossbox, below=0.5cm of scores, xshift=-0.6cm] (bpr) {$\mathcal{L}_{\text{BPR}}$};
\node[lossbox, below=0.5cm of scores, xshift=0.6cm] (coh) {$\lambda\!\cdot\!\mathcal{L}_{\text{coh}}$};

\draw[arr] (cust) -- (router);

\coordinate (fan-out) at ($(router.east) + (0.6,0)$);
\draw[line] (router.east) -- (fan-out);
\draw[arr] (fan-out) |- (sub0.west);
\draw[arr] (fan-out) |- (sub1.west);
\draw[arr] (fan-out) |- (sub2.west);
\draw[arr] (fan-out) |- (sub3.west);

\coordinate (fan-in) at ($(sub1.east)!0.5!(sub2.east) + (0.7,0)$);
\draw[line] (sub0.east) -- ++(0.7,0) |- (fan-in);
\draw[line] (sub1.east) -- ++(0.7,0) |- (fan-in);
\draw[line] (sub2.east) -- ++(0.7,0) |- (fan-in);
\draw[line] (sub3.east) -- ++(0.7,0) |- (fan-in);
\draw[arr] (fan-in) -- (scores.west);

\coordinate (train-split) at ($(scores.south) + (0,-0.2)$);
\draw[line] (scores.south) -- (train-split);
\draw[arr] (train-split) -| (bpr.north);
\draw[arr] (train-split) -| (coh.north);

\draw[arr] (scores.east) -- (topk.west);

\node[font=\tiny, below=0.02cm of topk.south, anchor=north] {Inference};
\node[font=\tiny, below=0.02cm of $(bpr.south)!0.5!(coh.south)$, anchor=north] {Training};

\begin{scope}[on background layer]
\node[groupbox, fit=(sub0)(sub3), label={[font=\scriptsize]above:Band-stratified LightGCN}] {};
\end{scope}

\end{tikzpicture}%
}
\caption{PC-LightGCN architecture. Each customer is routed by declared MiFID band to a dedicated sub-model. Training combines $\mathcal{L}_{\text{BPR}}$ with a coherence margin $\lambda \cdot \mathcal{L}_{\text{coh}}$. At inference the sub-model produces the top-$k$ list directly.}
\label{fig:architecture}
\end{figure*}
